\documentclass[12pt]{article}

\usepackage{fontawesome}
\usepackage{hyperref}
\usepackage{xurl}
\usepackage[tagged, highstructure]{accessibility}
\usepackage{bookmark}
\usepackage{enumitem}

\hypersetup{
    colorlinks=false,
    pdfborder={0 0 0},
    pdftitle={File Compression and Archives Lab},
    pdfauthor={Adrianna Holden-Gouveia},
    pdfsubject={Introduction to Linux - Lab Assignment},
    pdflang={en-US},
    pdfkeywords={compression, archives, tar, gzip}
}

% Configure PDF bookmarks for navigation
\bookmarksetup{
    numbered,
    open,
}

% Configure list spacing for better accessibility
\setlist{nosep}


\title{File Compression and Archives}
\author{
        Adrianna Holden-Gouveia \\
        Website: \href{https://aholdengouveia.name}{aholdengouveia.name}\\ 
        \date{\vspace{-5ex}}
        %Email: \href{mailto:admin@aholdengouveia.name}{admin@aholdengouveia.name} \\
%        \faLinkedin{: aholdengouveia} \\
        \faGithub {: aholdengouveia} \\
%        \faTwitter {: aholdengouveia} \\
        }

%S\date{\today}


\begin{document}    

\maketitle

\section*{Accessibility Notice}
This document is also available in HTML format at:

Link: \href{https://aholdengouveia.name/IntroLinux/labs/compressionarchive.html}{\textbf{\nolinkurl{https://aholdengouveia.name/IntroLinux/labs/compressionarchive.html}}}

The HTML version provides enhanced accessibility features including keyboard navigation, screen reader support, responsive design, dark mode support, and high contrast options.


%\begin{abstract}
%This is a template for Linux Administration Lab work
%\end{abstract}
%\tableofcontents

\section*{Learning Outcomes}
By the end of this lab, you will be able to:
\begin{itemize}
    \item Explain the difference between an archive and a compressed file
    \item Create, list, and extract tar archives
    \item Compress and decompress files using gzip
    \item Compare file sizes before and after compression
\end{itemize}

\section*{Lab Instructions}
In this lab you will bundle files into an archive, compress them, and compare the results. You can do this on any Linux system.

\subsection*{How to set up for this lab}
\begin{itemize}
    \item Create a directory for this lab, and put at least 5 files in it.
    \item Put some text in each file, so there is something to compress.
    \item Include at least one larger file. For example, copy /etc/services into your lab directory.
\end{itemize}

\subsection*{Problems to solve}
    \begin{enumerate}
        \item In your own words, what is the difference between an archive (such as a tar file) and a compressed file (such as a gzip file)?
        \item Archives
        \begin{enumerate}
            \item Use tar to create an archive of your lab directory. Take a screenshot.
            \item List the contents of your archive without extracting it. Take a screenshot.
            \item Make a new directory and extract your archive into it. Use ls to show the files are there. Take a screenshot.
        \end{enumerate}
        \item Compression
        \begin{enumerate}
            \item Use ls -l to check the size of your larger file.
            \item Compress it with gzip, then check the size again. How much smaller is it? Take a screenshot.
            \item Decompress it with gunzip. Take a screenshot.
        \end{enumerate}
        \item Archive and compress in one step
        \begin{enumerate}
            \item Use tar to create a compressed archive (.tar.gz) of your lab directory. Take a screenshot.
            \item Use ls -l to compare the sizes of your .tar and .tar.gz files. Which is smaller, and by how much? Take a screenshot.
            \item Extract your .tar.gz into a new directory. Take a screenshot.
        \end{enumerate}
        \item Give one example of when you would use an archive, one example of when you would use compression, and one example of when you would use both.
    \end{enumerate}

\section*{Deliverables}
One file with the lab title, your name, and the date at the top.
\begin{itemize}
    \item Each answer should be numbered to match the question.
    \item Make sure to list each command and the result.
    \item Make sure to take a screenshot of each successful command, labeled with which problem it solved. Your screenshots must show your prompt with your username, or a note with your name and the term.
\end{itemize}

\section*{Resources}
References, a video, a PowerPoint and some notes are available at my website
Link: \href{https://www.aholdengouveia.name/IntroLinux/compressionarchive.html}{\textbf{\nolinkurl{https://www.aholdengouveia.name/IntroLinux/compressionarchive.html}}}

My FAQ Linux playlist on YouTube, with demos of many of the commands and scripts in this class
Link: \href{https://www.youtube.com/playlist?list=PLQOmuLC45eoQlIMH7oCNL6yFP7kWQubBo}{\textbf{\nolinkurl{https://www.youtube.com/playlist?list=PLQOmuLC45eoQlIMH7oCNL6yFP7kWQubBo}}}

My video on how to use tar
Link: \href{https://www.youtube.com/watch?v=x0EK6PmKRCU}{\textbf{\nolinkurl{https://www.youtube.com/watch?v=x0EK6PmKRCU}}}

\end{document}